% !TEX root = ../main.tex
\section{Respuestas a las preguntas de análisis}
% La guía solicita responder cada pregunta; se conserva el formato de consulta breve.
\begin{enumerate}
\setlength{\itemsep}{.2em}
\item \textbf{¿Por qué basta probar divisores hasta $\sqrt n$?} Si $n=ab$ y $a>\sqrt n$, $b>\sqrt n$, entonces $ab>n$, contradicción. Todo compuesto tiene al menos un divisor $\le\sqrt n$.

\item \textbf{¿Por qué esa cota es exponencial en bits?} Si $n$ tiene $L$ bits, $\sqrt n$ es del orden de $2^{L/2}$: exponencial en $L$, aunque sublineal en $n$.

\item \textbf{¿Por qué Fermat puede demostrar composición, pero no primalidad?} Para un primo $p$ y base coprima $a$, Fermat exige $a^{p-1}\equiv1\pmod p$. Si falla, por contraposición $n$ es compuesto; si pasa, puede ser un pseudoprimo. El ejemplo $341$ con base $2$ confirma que el recíproco es falso.

\item \textbf{¿Qué diferencia hay entre pseudoprimo y número de Carmichael?} Un pseudoprimo pasa Fermat para alguna base coprima; un Carmichael, para todas. $341$ pasa base $2$; $561$ pasa cualquier base coprima, pero la base $3$ revela un factor.

\item \textbf{¿Por qué Miller--Rabin examina $n-1=2^sd$?} Separar la máxima potencia de $2$ produce la cadena $a^d,a^{2d},\ldots,a^{2^{s-1}d}$. En un módulo primo, si el último término es $1$, la cadena debe comenzar en $1$ o alcanzar $-1$ antes: solo $\pm1$ son raíces de $x^2=1$ en un cuerpo.

\item \textbf{¿Qué significa encontrar una raíz no trivial de $1$?} Si $x\not\equiv\pm1\pmod n$ y $x^2\equiv1\pmod n$, entonces $n$ es compuesto. Para $341$, $x=32$ y $\gcd(x\pm1,341)$ recupera $31$ y $11$.

\pagebreak[4]
\item \textbf{¿Qué diferencia hay entre testigo y mentiroso fuerte?} El testigo rechaza un compuesto; el mentiroso fuerte lo acepta. Para $341$, la base $2$ pasa Fermat, pero es testigo de Miller--Rabin.

\item \textbf{¿Por qué el error es unilateral?} Toda base válida hace pasar a un número primo. El rechazo significa compuesto con certeza; la única posible equivocación es aceptar como probable primo un compuesto que superó las bases elegidas.

\item \textbf{¿Qué hipótesis permite usar $4^{-r}$?} El número evaluado debe ser compuesto impar; cada base se selecciona uniforme e independientemente entre las bases válidas. La fracción de mentirosos fuertes es $\le1/4$ y el evento de $r$ aceptaciones tiene probabilidad $\le(1/4)^r$. La semilla fija permite reproducir un experimento, no garantiza aleatoriedad criptográfica.

\item \textbf{¿Por qué una semilla fija no es apropiada para generar claves?} Permite reconstruir candidatos y bases si se conoce o adivina; se requiere una fuente impredecible como \texttt{secrets.SystemRandom}.

\item \textbf{¿Cómo influye el precribado?} Descarta múltiplos de primos pequeños y certifica si el candidato coincide con uno de ellos. Ahorra exponenciaciones, sin certificar números grandes que lo superan.

\item \textbf{¿Por qué ``no observamos errores'' no significa ``el algoritmo no puede errar''?} Una muestra finita puede omitir las bases mentirosas. La ausencia observada no prueba probabilidad cero; hay que distinguir frecuencia empírica, probabilidad exacta $p_n^r$ para un compuesto dado y cota universal $4^{-r}$.
\end{enumerate}
